\documentclass[../DoAn.tex]{subfiles}
\begin{document}

Chương 4 đi sâu vào bước bản lề quyết định độ chính xác của mô hình học sâu: phân tích và xử lý đặc trưng lưu lượng mạng. Điểm nhấn của chương là việc đánh giá, so sánh các bộ dữ liệu tiêu chuẩn dùng trong NIDS. Trên cơ sở đó, chương giải thích quyết định lựa chọn CIC-IDS-2017 làm bộ dữ liệu nghiên cứu cốt lõi, thay thế cho NSL-KDD truyền thống. Tiếp theo, chương trình bày các phương pháp tiền xử lý dữ liệu và thiết kế luồng tích hợp hệ thống phần mềm nhằm vận hành mô hình trong điều kiện thời gian thực.

\section{Phân tích các bộ dữ liệu đánh giá NIDS}

\subsection{Tập dữ liệu NSL-KDD (Bộ dữ liệu nền tảng)}
NSL-KDD là phiên bản tinh chỉnh và nâng cấp của bộ dữ liệu KDD Cup 99 nổi tiếng. Bằng cách loại bỏ các bản ghi trùng lặp (redundant records), NSL-KDD giúp các thuật toán học máy không bị thiên lệch về phía dữ liệu dư thừa và cung cấp kết quả đánh giá mô hình khách quan hơn. Bộ dữ liệu này biểu diễn mỗi luồng kết nối bằng 41 đặc trưng (bao gồm đặc trưng cơ bản, đặc trưng nội dung và đặc trưng thống kê lưu lượng) và phân thành 4 nhóm tấn công chính: DoS (Denial of Service), Probe (Quét mạng), R2L (Remote to Local) và U2R (User to Root).

Trong nhiều năm qua, NSL-KDD được xem là "bài kiểm tra chuẩn mực" (benchmark) cho các nghiên cứu về IDS và đã được sử dụng rộng rãi trong các đồ án tiền nhiệm để chứng minh tính khả thi ban đầu của các mô hình AI. Tuy nhiên, giới hạn lớn nhất của bộ dữ liệu này là tuổi đời: dữ liệu gốc được thu thập từ năm 1999, do đó không phản ánh được cấu trúc liên kết mạng hiện đại cũng như các thủ đoạn tấn công mới xuất hiện trong thập kỷ qua.

\subsection{Tập dữ liệu CIC-IDS-2017 (Bộ dữ liệu thực chiến chính)}
Để bắt kịp với sự tiến hóa của các mối đe dọa không gian mạng, Viện An ninh Thông tin Canada (CIC) đã phát hành bộ dữ liệu CIC-IDS-2017. Bộ dữ liệu này được xây dựng trong một môi trường giả lập tiệm cận thực tế, thu thập dữ liệu thô (PCAP) trong 5 ngày với sự tham gia của nhiều thiết bị hiện đại.

Điểm nổi bật của CIC-IDS-2017 là việc sử dụng công cụ CICFlowMeter để tự động trích xuất hơn 78 đặc trưng thống kê luồng (Flow-based features) như thời gian tồn tại của luồng, tốc độ truyền gói tin, hay kích thước cửa sổ TCP. Bộ dữ liệu chứa hơn 2.8 triệu bản ghi, bao trùm luồng dữ liệu an toàn (Benign) và 14 loại tấn công hiện đại nhất như: Brute Force (FTP/SSH), DoS/DDoS, Heartbleed, Web Attack (SQL Injection, XSS), Infiltration và Botnet.

\subsection{So sánh và lý do chuyển đổi sang CIC-IDS-2017}

\begin{table}[htbp]
\centering
\caption{So sánh đặc tính giữa bộ dữ liệu NSL-KDD và CIC-IDS-2017}
\label{tab:dataset_comparison}
\begin{tabular}{|p{3.5cm}|p{5cm}|p{5.5cm}|}
\hline
\textbf{Đặc tính} & \textbf{NSL-KDD} & \textbf{CIC-IDS-2017} \\
\hline
\textbf{Năm phát hành} & 1999 (KDD Cup) / 2009 & 2017 \\
\hline
\textbf{Đặc trưng} & 41 đặc trưng (dựa trên nội dung/kết nối) & 78+ đặc trưng (dựa trên hành vi luồng - Flow-based) \\
\hline
\textbf{Các loại tấn công} & DoS, Probe, R2L, U2R & Brute Force, DoS/DDoS, Web Attack, Infiltration, Botnet \\
\hline
\textbf{Tính đại diện} & Thấp (Lỗi thời, thiếu giao thức mới) & Rất cao (Mô phỏng mạng hiện đại, có mã hóa) \\
\hline
\textbf{Mất cân bằng lớp} & Đã được cân bằng nhân tạo & Mất cân bằng tự nhiên (như thực tế) \\
\hline
\end{tabular}
\end{table}

Dù NSL-KDD đóng vai trò quan trọng trong việc thử nghiệm các thuật toán cơ bản, đồ án này quyết định chuyển hướng tập trung hoàn toàn vào CIC-IDS-2017 làm bộ dữ liệu huấn luyện và đánh giá chính. Sự chuyển đổi này xuất phát từ ba lý do cốt lõi:

\begin{enumerate}
    \item \textbf{Tính thời sự và tính đa dạng của rủi ro (Threat Landscape)}: NSL-KDD hoàn toàn thiếu vắng các hình thức tấn công nguy hiểm thời nay như Web Attack (tấn công ứng dụng web), Botnet hay Infiltration (xâm nhập nội bộ). CIC-IDS-2017 bao phủ đầy đủ các rủi ro này, giúp mô hình AI của đồ án có giá trị phòng thủ thực tế.
    \item \textbf{Định dạng đặc trưng (Feature Representation)}: Đặc trưng của NSL-KDD dựa trên việc trích xuất nội dung (payload) vốn tiêu tốn nhiều thời gian và gặp khó khăn khi mã hóa (encryption) lên ngôi. Ngược lại, 78+ đặc trưng của CIC-IDS-2017 hoàn toàn dựa trên hành vi của dòng chảy mạng (Flow-based) sinh ra từ CICFlowMeter. Sự tương đồng này cho phép đồ án tích hợp trực tiếp CICFlowMeter vào Data Pipeline thời gian thực, đảm bảo tính đồng nhất giữa lúc huấn luyện và khi chạy thực chiến (inference).
    \item \textbf{Thách thức về mất cân bằng lớp (Class Imbalance)}: CIC-IDS-2017 phản ánh chính xác đặc thù của mạng lưới thực: lưu lượng sạch chiếm áp đảo (hơn 80\%), trong khi một số tấn công như Infiltration hay Heartbleed chiếm chưa tới 0.1\%. Bài toán mất cân bằng cực đoan này là môi trường thử thách hoàn hảo để chứng minh sức mạnh của mô hình FT-Transformer và cơ chế Hybrid Ensemble mà đồ án đề xuất, điều mà NSL-KDD không thể làm được.
\end{enumerate}

\begin{table}[htbp]
\centering
\caption{Phân phối các lớp tấn công tiêu biểu trong bộ dữ liệu CIC-IDS-2017}
\label{tab:cic_distribution}
\begin{tabular}{|l|r|c|}
\hline
\textbf{Loại lưu lượng (Label)} & \textbf{Số lượng mẫu ước tính} & \textbf{Tỷ lệ (\%)} \\
\hline
Benign (Bình thường) & 2.271.320 & 80.30\% \\
\hline
DoS / DDoS & 380.699 & 13.46\% \\
\hline
PortScan & 158.930 & 5.62\% \\
\hline
Brute Force (FTP/SSH) & 13.835 & 0.49\% \\
\hline
Web Attack & 2.180 & 0.08\% \\
\hline
Botnet & 1.966 & 0.07\% \\
\hline
Infiltration & 36 & $<$ 0.01\% \\
\hline
\end{tabular}
\end{table}

\section{Tiền xử lý dữ liệu và Kỹ thuật Feature Engineering}
Để chuẩn bị cho mô hình học sâu, dữ liệu thô từ CIC-IDS-2017 phải trải qua các bước tiền xử lý nghiêm ngặt:
\begin{itemize}
    \item \textbf{Làm sạch dữ liệu (Data Cleaning)}: Loại bỏ các bản ghi chứa giá trị rỗng (NaN) hoặc vô cực (Infinity) sinh ra trong quá trình trích xuất của CICFlowMeter.
    \item \textbf{Chuẩn hóa (Normalization/Scaling)}: Do sự khác biệt cực lớn về đơn vị đo lường (ví dụ: kích thước tính bằng Byte, trong khi tốc độ tính bằng Mbps), các kỹ thuật Min-Max Scaling hoặc Robust Scaler được áp dụng để đưa dữ liệu về cùng một phổ tỷ lệ, giúp mạng nơ-ron hội tụ nhanh hơn.
    \item \textbf{Xử lý mất cân bằng (Oversampling/Undersampling)}: Đối với các tấn công hiếm như Botnet, kỹ thuật sinh dữ liệu nhân tạo (SMOTE) có thể được kết hợp ở mức độ vừa phải để hỗ trợ nhóm mô hình Tree-based trong kiến trúc Ensemble tránh bị thiên lệch.
\end{itemize}

\section{Quá trình phát triển và Phân tích thuật toán Học sâu}
Để đạt được độ chính xác tối đa trong môi trường thực chiến, hệ thống học sâu của đồ án đã trải qua nhiều giai đoạn phát triển và tối ưu hóa thuật toán. Việc đánh giá sự tiến hóa của các phiên bản mô hình giúp làm rõ lý do tại sao kiến trúc cuối cùng được lựa chọn.

\subsection{Phiên bản 1 (Cơ sở): Mạng nơ-ron nhân tạo (ANN) và FT-Transformer V1}
Trong giai đoạn đầu của nghiên cứu, đồ án triển khai một Mạng nơ-ron nhân tạo (ANN - Artificial Neural Network) truyền thống làm mô hình cơ sở (baseline). ANN bao gồm các lớp ẩn kết nối đầy đủ (fully-connected) đi kèm hàm kích hoạt ReLU. Dù ANN có tốc độ suy luận nhanh và dễ triển khai, nó tỏ ra thiếu hiệu quả trong việc nắm bắt các mối quan hệ phức tạp giữa 78 đặc trưng của luồng mạng, dẫn đến tỷ lệ phát hiện các cuộc tấn công hiếm (như Infiltration) rất thấp.

Nhận thấy hạn chế đó, phiên bản tiếp theo ứng dụng kiến trúc FT-Transformer (Feature Tokenizer Transformer) thế hệ đầu tiên. Khác với ANN, FT-Transformer biến từng đặc trưng số học thành một "token" (tương tự như từ vựng trong xử lý ngôn ngữ tự nhiên), sau đó dùng cơ chế Tự chú ý (Multi-Head Self-Attention) để đánh giá tầm quan trọng của từng đặc trưng. Phiên bản V1 này cải thiện rõ rệt khả năng nhận diện các dạng tấn công phức tạp. Tuy nhiên, khi đối mặt với bộ dữ liệu khổng lồ và nhiễu của CIC-IDS-2017, mạng Transformer có xu hướng bị học vẹt (overfitting) trên tập dữ liệu đa số (Benign) và thiếu ổn định khi tăng độ sâu của mạng.

\subsection{Phiên bản 2 (Tối ưu): FT-Transformer V2 và Hybrid Ensemble}
Để khắc phục bài toán overfitting và gia tăng độ ổn định, đồ án đã nâng cấp lên \textbf{FT-Transformer V2}. Quá trình xử lý và các cải tiến đột phá được mô tả chi tiết qua các hàm toán học sau:

\textbf{1. Hàm Nhúng đặc trưng (Feature Tokenizer)}
\begin{itemize}
    \item \textbf{Mục đích}: Biến đổi từng giá trị đặc trưng vô hướng (scalar) của CICFlowMeter thành một vector có số chiều không đổi ($d$) để phù hợp với đầu vào của mạng Transformer.
    \item \textbf{Công thức}:
    \begin{equation}
        T_j^{(0)} = W_j x_j + b_j
    \end{equation}
    Trong đó, $x_j$ là đặc trưng thứ $j$ của dòng chảy mạng, $W_j$ và $b_j$ là các ma trận trọng số và độ lệch (bias) có thể học được trong quá trình huấn luyện.
\end{itemize}

\textbf{2. Cơ chế Tự chú ý đa đầu (Multi-Head Self-Attention - MHSA)}
\begin{itemize}
    \item \textbf{Mục đích}: Đây là trái tim của mạng, cho phép đánh giá mối tương quan (tầm quan trọng chéo) giữa các đặc trưng khác nhau của một gói tin (ví dụ tương quan giữa thời lượng luồng và tốc độ truyền).
    \item \textbf{Công thức}:
    \begin{equation}
        \text{Attention}(Q, K, V) = \text{softmax}\left(\frac{Q K^T}{\sqrt{d_k}}\right) V
    \end{equation}
    Trong đó $Q$ (Query), $K$ (Key), $V$ (Value) là các ma trận biểu diễn từ token đầu vào. Việc chia cho $\sqrt{d_k}$ giúp tránh hiện tượng gradient bùng nổ khi tính hàm softmax.
\end{itemize}

\textbf{3. Hàm LayerScale}
\begin{itemize}
    \item \textbf{Mục đích}: Cung cấp các tham số có thể học được để tỷ lệ hóa (scale) đầu ra của từng khối Attention trước khi cộng dồn (residual). Cơ chế này giúp luồng đạo hàm truyền qua các mạng sâu ổn định hơn.
    \item \textbf{Công thức}:
    \begin{equation}
        x_{l+1} = x_l + \text{diag}(\lambda_{l,1}, \dots, \lambda_{l,d}) \cdot \text{MHSA}(x_l)
    \end{equation}
    Trong đó $\lambda_{l,i}$ là siêu tham số học được khởi tạo cực nhỏ (vd: $10^{-4}$), giúp mạng ban đầu hoạt động giống mạng nông sau đó mới tối ưu hóa dần.
\end{itemize}

\textbf{4. Hàm DropPath (Stochastic Depth)}
\begin{itemize}
    \item \textbf{Mục đích}: Hoạt động như một cơ chế chính quy hóa (regularization) triệt để. Thay vì chỉ Dropout node nơ-ron cục bộ, DropPath bỏ qua ngẫu nhiên toàn bộ một khối mạng khi huấn luyện, ép mô hình không được học vẹt (overfitting).
    \item \textbf{Công thức}:
    \begin{equation}
        H_{l+1} = x_l + p_l \cdot F(x_l)
    \end{equation}
    Trong đó $p_l \sim \text{Bernoulli}(1 - p_{\text{drop}})$ là biến ngẫu nhiên nhận giá trị $0$ hoặc $1$. Nếu $p_l = 0$, khối $F(x_l)$ hoàn toàn bị tắt.
\end{itemize}

Bên cạnh đó, đồ án thiết kế cơ chế \textbf{Hybrid Ensemble} (Hội đồng biểu quyết). Mô hình Transformer được kết hợp cùng Rừng ngẫu nhiên (Random Forest) và K-Nearest Neighbors (KNN). Cơ chế Soft Voting được sử dụng với công thức:
\begin{equation}
    \hat{P} = \alpha \cdot P_{\text{Transformer}} + \beta \cdot P_{\text{RF}} + \gamma \cdot P_{\text{KNN}}
\end{equation}
Trong đó $\alpha, \beta, \gamma$ là các trọng số tinh chỉnh dựa trên độ tin cậy của từng mô hình. Điều này bù đắp khả năng xử lý phân phối dữ liệu phi tuyến tính cục bộ mà Transformer bỏ sót.

\subsection{So sánh sự khác biệt giữa các phiên bản mô hình}

\begin{table}[htbp]
\centering
\caption{So sánh sự tiến hóa giữa các phiên bản thuật toán trong đồ án}
\label{tab:model_comparison}
\begin{tabular}{|p{3.5cm}|p{4.5cm}|p{6cm}|}
\hline
\textbf{Tiêu chí} & \textbf{Phiên bản 1 (ANN \& FT-Transformer V1)} & \textbf{Phiên bản 2 (FT-Transformer V2 \& Hybrid Ensemble)} \\
\hline
\textbf{Khả năng trích xuất đặc trưng} & Tuyến tính (ANN) hoặc chú ý cục bộ, dễ bị nhiễu (FT-Transformer V1). & Tối ưu hóa sâu nhờ DropPath và LayerScale, tập trung vào đặc trưng lõi. \\
\hline
\textbf{Hiện tượng học vẹt (Overfitting)} & Dễ xảy ra do mất cân bằng dữ liệu của tập CIC-IDS-2017. & Được kiểm soát chặt chẽ nhờ Stochastic Depth (DropPath). \\
\hline
\textbf{Nhận diện tấn công hiếm (Infiltration)} & Kém (Thường gán nhầm thành Benign do mất cân bằng). & Tốt (Nhờ sự bù trừ của RF và KNN trong cơ chế Ensemble). \\
\hline
\textbf{Tỷ lệ cảnh báo giả (FPR)} & Khá cao (do mô hình đơn lẻ thiếu sự kiềm chế chéo). & Cực kỳ thấp (triệt tiêu sai số thông qua cơ chế biểu quyết). \\
\hline
\end{tabular}
\end{table}

Việc chuyển đổi từ một mô hình học sâu cơ bản sang kiến trúc kết hợp (Hybrid Ensemble) với lõi FT-Transformer V2 minh chứng cho một bước tiến dài trong việc giải quyết bài toán hóc búa của phát hiện xâm nhập mạng.

\section{Thiết kế luồng dữ liệu và Tích hợp hệ thống phân tích thời gian thực}

\vspace{1em}
\textbf{Kết chương}
Chương 4 đã làm rõ tầm quan trọng của việc lựa chọn dữ liệu thông qua phép so sánh trực diện giữa NSL-KDD và CIC-IDS-2017. Bằng cách sử dụng CIC-IDS-2017, hệ thống AI của đồ án được định hình để đối mặt với những bài toán hóc búa nhất của an ninh mạng hiện đại. Đồng thời, quy trình tiền xử lý được thiết kế chặt chẽ không chỉ phục vụ huấn luyện mà còn tích hợp liền mạch vào kiến trúc thời gian thực của ứng dụng. Mức độ hiệu quả của việc phân tích này sẽ được chứng minh thông qua các kết quả thực nghiệm trình bày ở Chương 5.

\end{document}
